\documentclass[11pt,a4paper,landscape]{article}

% Pakete
\usepackage[utf8]{inputenc}
\usepackage[T1]{fontenc}
\usepackage[ngerman]{babel}
\usepackage[left=15mm,right=15mm,top=12mm,bottom=12mm]{geometry}
\usepackage{amsmath,amssymb}
\usepackage{tikz}
\usepackage{siunitx}
\sisetup{locale=DE, per-mode=fraction}

% Serifenlose Schrift
\usepackage{helvet}
\renewcommand{\familydefault}{\sfdefault}

% Keine Seitenzahl
\pagestyle{empty}

% Keine Einrückung
\setlength{\parindent}{0pt}
\setlength{\parskip}{0pt}

\begin{document}

% ============================================================
% TITEL
% ============================================================

\begin{center}
{\LARGE\bfseries De-Broglie-Wellenlänge}\\[0.3em]
\rule{0.5\textwidth}{1.5pt}
\end{center}

\vspace{0.5em}

% ============================================================
% ZWEI SPALTEN
% ============================================================

\begin{minipage}[t]{0.48\textwidth}

% ============================================================
% KONTEXT
% ============================================================

\fbox{\parbox{0.95\textwidth}{
\begin{center}
{\large\bfseries Kontext: Von Einstein zu De Broglie}
\end{center}
}}

\vspace{0.5em}

\begin{itemize}
    \setlength{\itemsep}{4pt}
    \item \textbf{1905 Einstein:} Licht zeigt Teilcheneigenschaften (Photoeffekt)
    \item \textbf{1924 De Broglie:} Umkehrschluss -- haben Teilchen auch Welleneigenschaften?
    \item \textbf{1927 Davisson-Germer:} Experimenteller Nachweis durch Elektronenbeugung an Kristallen
\end{itemize}

\vspace{0.8em}

% ============================================================
% FORMEL
% ============================================================

\fbox{\parbox{0.95\textwidth}{
\begin{center}
{\large\bfseries De-Broglie-Hypothese}
\end{center}
}}

\vspace{0.5em}

\textit{,,Jedes bewegte Teilchen hat eine Wellenlänge:''}

\vspace{0.3em}
\begin{center}
\fbox{\fbox{\parbox{0.7\textwidth}{
\begin{center}
{\Large $\lambda = \dfrac{h}{p} = \dfrac{h}{m \cdot v}$}
\end{center}
}}}
\end{center}

\vspace{0.5em}

\textbf{Für beschleunigte Elektronen:}

\vspace{0.2em}
\begin{tabular}{@{}r@{\;$\Rightarrow$\;}l@{\quad}l@{}}
$e \cdot U = \frac{1}{2} m_e v^2$ & $v = \sqrt{\dfrac{2 \cdot e \cdot U}{m_e}}$ & {\small(Energieerhaltung)} \\[0.8em]
$\lambda = \dfrac{h}{m_e \cdot v}$ & $\lambda = \dfrac{h}{\sqrt{2 \cdot m_e \cdot e \cdot U}}$ & {\small(einsetzen)}
\end{tabular}

\end{minipage}
\hfill
\begin{minipage}[t]{0.48\textwidth}

% ============================================================
% GRÖSSENORDNUNGEN
% ============================================================

\fbox{\parbox{0.95\textwidth}{
\begin{center}
{\large\bfseries Größenordnungen}
\end{center}
}}

\vspace{0.5em}

\begin{tabular}{@{}l@{\quad}r@{\quad}l@{}}
Elektron (50 V): & $\lambda \approx$ & $\SI{170}{pm}$ \quad {\small(Atomgröße!)} \\[0.3em]
Elektron (50 kV): & $\lambda \approx$ & $\SI{5}{pm}$ \\[0.3em]
Fußball (30 m/s): & $\lambda \approx$ & $\SI{e-35}{m}$ \quad {\small(nicht messbar)}
\end{tabular}

\vspace{0.3em}

$\Rightarrow$ Je größer $m$ oder $v$, desto kleiner $\lambda$

\vspace{0.8em}

% ============================================================
% BEDEUTUNG
% ============================================================

\fbox{\parbox{0.95\textwidth}{
\begin{center}
{\large\bfseries Bedeutung \& Anwendung}
\end{center}
}}

\vspace{0.5em}

\textbf{Welle-Teilchen-Dualismus ist universal:}\\
Nicht nur Licht -- \textit{alle} Materie hat Wellen- und Teilcheneigenschaften!

\vspace{0.5em}

\textbf{Anwendung: Elektronenmikroskop}

\vspace{0.3em}
\fbox{\parbox{0.9\textwidth}{
\textbf{Auflösung begrenzt durch Wellenlänge} (Beugungslimit):\\[0.3em]
\begin{tabular}{@{}ll@{}}
Lichtmikroskop: & $\lambda \approx \SI{500}{nm}$ $\Rightarrow$ max.\ $\sim\SI{200}{nm}$ \\[0.2em]
Elektronenmikroskop: & $\lambda \approx \SI{5}{pm}$ $\Rightarrow$ \textbf{Atome sichtbar!}
\end{tabular}
}}

\vspace{0.4em}

{\small $\Rightarrow$ Elektronen haben bei hoher Spannung viel kleinere $\lambda$ als Licht\\
\hspace*{0.7em} $\Rightarrow$ Feinere Strukturen auflösbar}

\end{minipage}

\vspace{0.8em}

% ============================================================
% KONSTANTEN
% ============================================================

\begin{center}
\rule{0.9\textwidth}{0.5pt}
\vspace{0.3em}

\textbf{Konstanten:}\quad
$h = \SI{6,63e-34}{J.s}$
\hspace{1.5cm}
$e = \SI{1,6e-19}{C}$
\hspace{1.5cm}
$m_e = \SI{9,1e-31}{kg}$
\hspace{1.5cm}
$\SI{1}{pm} = \SI{e-12}{m}$
\end{center}

\end{document}
